\documentclass[lettersize,journal]{IEEEtran}
\usepackage{amsmath,amsfonts}
\usepackage{algorithmic}
\usepackage{algorithm}
\usepackage{array}
\usepackage[caption=false,font=normalsize,labelfont=sf,textfont=sf]{subfig}
\usepackage{textcomp}
%\usepackage{stfloats}
\usepackage{url}
\usepackage{verbatim}
\usepackage{graphicx}
\usepackage{cite}

\begin{document}

\title{Titolo}

\author{Nome Cognome Matricola,}

\maketitle

\begin{abstract}
L'abstract serve a riassumere brevemente quello che presenta il report. Per istruzioni dettagliate sul formato fare riferimento al file ``New\_IEEEtran\_how-to.tex''.
\end{abstract}

\section{Introduzione}
Introdurre il problema e presentare brevemente i dati.

\section{Analisi dei dati}
Caratterizzare i dati tramite grafici e descrizione testuale.

\section{Metodologia}
Indicare chiaramente le scelte effettuate per i modelli e una breve descrizione del loro funzionamento e del perchè sono stati usati. 

\section{Risultati sperimentali}
Presentare i risultati ottenuti mettendo in luce forze e debolezze delle varie proposte. Se ritenuto utile riportare anche degli esempi pratici per faciliatere la comprensione. Indicare i parametri usati per i risultati presentati.

\section{Conclusioni}
Indicare come è possibile sfruttare i dati per un applicazione di business.

\section{Contributi}
Indicare i contributi personali.

\end{document}


